\documentclass{article}
\usepackage[utf8]{inputenc}
\usepackage[french]{babel}
\usepackage{amsmath}
\usepackage{graphicx}
\usepackage{listings}
\usepackage{xcolor}
\usepackage{comment}
\usepackage[top=2cm, bottom=2cm, left=2cm, right=2cm]{geometry}

% Configuration pour le code Python
\lstset{
    language=Python,
    basicstyle=\ttfamily\small,
    keywordstyle=\color{blue},
    commentstyle=\color{gray},
    stringstyle=\color{red},
    showstringspaces=false,
    breaklines=true,
    frame=single,
    numbers=left,
    numberstyle=\tiny\color{gray}
}

\title{Complexité et notation \emph{Big O}}
\date{Novembre 2025}

\begin{document}

\maketitle

\subsubsection*{Ce document couvre les notions de complexité des algorithmes et en particulier la notation \emph{big O}}

La notation \emph{Big O} a pour but d'exprimer la complexité temporelle des algorithmes dans le pire scenario d'exécution (\textit{Worst Case scenario}) et suit quelques principes de base:

\begin{itemize}
    \item Les constantes sont ignorées
    
    $O(2n) = 2 \cdot O(n) = O(n)$
    
    \item Les termes dominés sont ignorés
    
    $O(2n^2 + 5n + 50) = O(n^2)$
    
    % \item Utilisez différentes variables pour différents inputs
    
    \item Uniquement le pire cas d'exécution est considéré
\end{itemize}

\section*{Complexités les plus fréquentes}

Les complexités les plus fréquentes sont les suivantes:

\subsubsection*{Constante $O(1)$}

Cela veut dire que la complexité ne depend pas de la taille du vecteur en input.
\begin{itemize}
    \item Les méthodes \emph{append} et \emph{pop} sur les listes en Python.
\end{itemize}

\subsubsection*{Logarithmique $O(\log(n))$}

\begin{itemize}
    \item L'algorithme de recherche binaire
\end{itemize}

\subsubsection*{Linéaire $O(n)$}

\begin{itemize}
    \item Les boucles \emph{for} et \emph{while}
\end{itemize}

\subsubsection*{Quadratique, Cubique, etc. $O(n^2)$, $O(n^3)$}

\begin{itemize}
    \item Boucles imbriquées ayant le même nombre d'itérations
    \item L'algorithme de tri par insertion tel que vu en cours à une complexité de $O(n^2)$
\end{itemize}

\subsubsection*{Exponentielle $O(2^n)$}

\begin{itemize}
    \item Fibonacci par récursivité comme vu en exercice
\end{itemize}

\begin{comment}
\begin{figure}[h]
    \centering
    \includegraphics[width=0.6\textwidth]{complexity_chart.png}
    \caption{Complexité}
    \label{fig:complexity}
\end{figure}
\end{comment}

\section*{Exemples}

La boucle ci-dessous va exécuter \emph{n} fois (le \emph{n}  ici est choisi arbitrairement et représente la taille du vecteur \emph{array}) une opération de complexité $O(1)$, la complexité totale est donc $n \cdot O(1) = O(n)$.

\begin{lstlisting}
for number in array:
    array.pop()
\end{lstlisting}

Dans le prochain exemple nous avons deux vecteurs, si leur taille est respectivement de \emph{n} et \emph{m} alors la complexité sera de $O(n \cdot m)$ et non $O(n^2)$. Cela s'explique par le fait que les deux boucles imbriquées n'itèrent pas sur le même vecteur.

\begin{lstlisting}
def intersection_size(arrayA,arrayB):
    count = 0
    for a in arrayA:
        for b in arrayB:
            if a == b:
                count += 1
    return count
\end{lstlisting}

Les deux fonctions ci-dessous ont la même complexité de $O(n)$

\begin{lstlisting}
def find_extrem_1(array):
    mini , maxi = (array[0], )*2   
    for numbers in array:          
        if numbers < mini:
            mini = numbers
    for numbers in array:          
        if numbers > maxi: 
            maxi = numbers
    return mini , maxi

def find_extrem_2(array):
    mini , maxi = (array[0], )*2   
    for numbers in array:          
        if numbers < mini:
            mini = numbers
        elif numbers > maxi:
            maxi = numbers
    return mini , maxi
\end{lstlisting}

\end{document}